\scanpage{133}
3). La preuve de la formule (6) est la plus difficile.
Nous commençons par donner une
preuve directe et très simple pour le cas où
\deleted{$\unsure{p,q,r\geq1,\ r=1}$} \inserted{$p\geq1,q\geq1,r=1$} *
(preuve qui a d'ailleurs l'avantage
de \inserted{devoir} s'appliquer sans modification au cas d'une
trace \unsure{$\tau$|$\sigma$|$\operatorname{Tr}$|$D$} au sens de von Neumann - Dixmier, définie
sur une algèbre auto-adjointe $\unsure{\mathcal A|\mathfrak A}$ \deleted{\unsure{ft fermée}}
\deleted{\unsure{unift fermée}} unif\textsuperscript{t} fermée \unsure{fortement}
(1) (\deleted{\unsure{les éléments de}} \textsuperscript{(1)} les ``\deleted{\unsure{éléments}} \inserted{opérateurs} de puissance
$p$\textsuperscript{ième} sommable'' de $\unsure{\mathcal A|\mathfrak A}$ \inserted{-- au sens de la trace $\unsure{\tau|\sigma|\operatorname{Tr}|D}$ --}
se définissant encore
comme l'ens des $A\in\unsure{\mathcal A|\mathfrak A}$ tels que
$\unsure{\tau|\sigma|\operatorname{Tr}|D}(|A|^p)<+\infty$.)\textsuperscript{\unsure{*|a|A}}.
\margin{?? Il n'est d'ailleurs pas difficile \deleted{\unsure{de}} étendre alors le même \deleted{\unsure{tout}} énoncé au cas où $r\geq1$}
(Au N° suivant, nous
obtiendrons d'ailleurs le $=$ résultat comme
conséquence d'une proposition plus générale).

Nous commençons par prouver que, si
$A_1,\ldots A_{2^n}$ sont $2^n$ opérateurs cpt continus,
($n$ entier $\geq1$) alors on a \textsuperscript{(1)}
\[
\|A_1 A_2\mathinner{\cdotp\cdotp\cdotp\cdotp} A_{2^n}\|_1\leq\|A_1\|_{2^n}\cdot\|A_2\|_{2^n}\ldots\|A_{2^n}\|_{2^n}\tag*{(11)}
\]
Nous le prouvons par récurrence sur $n$\,;
c'est vrai si $n=1$, \deleted{\unsure{on a bien}} \inserted{(formule\quad)}, supposons
\deleted{\unsure{$\|A_1A_2\|_1\leq\|A_1\|_2$}}
-la démontrée \inserted{\unsure{la|le} formule analogue} pour $2^{n-1}$ facteurs, et démontrons
alors (11). Écrivant $A_1\ldots A_{2^n}=(A_1A_2)\ldots(A_{2^n-1}A_{2^n})$
le premier membre sera, en vertu de l'hyp.
de récurrence, majoré par $\|A_1A_2\|_{2^{n-1}}\ldots\|A_{2^n-1}A_{2^n}\|_{2^{\unsure{n-1}}}$

\noindent\rule{\linewidth}{0.3pt}

(1) Le lecteur notera que ce résultat suffirait pour
démontrer \deleted{\unsure{directement|exactement}} \inserted{directement}, comme nous le ferons plus bas,
la formule \unsure{(6)|(6')} dans le cas général.
